\chapter{Defibrillation}

\section*{What Is This Test or Procedure?}
Defibrillation delivers a therapeutic dose of electrical energy to the heart with a device called a defibrillator. This depolarizes the heart muscle, terminating dysrhythmias and allowing the SA node to re-establish a normal rhythm.

\section*{Alternative Names}
Electrical shock therapy, Cardioversion

\section*{Principle}
A powerful, unsynchronized electrical shock depolarizes the entire myocardium at once, terminating the disorganized activity of ventricular fibrillation or pulseless ventricular tachycardia. The resulting pause allows the heart's normal pacemaker to re-establish a coordinated rhythm.
\section*{History}
Carl Ludwig and Florian Prevost demonstrated electrical termination of VF in dogs in 1899. Claude Beck performed the first successful human defibrillation during surgery in 1947.

\section*{Why This Test or Procedure Is Ordered}
Ordered immediately for pulseless Ventricular Fibrillation (VF) or pulseless Ventricular Tachycardia (VT) during cardiac arrest.

\section*{Is This a Primary or Secondary Test or Procedure}
Primary curative treatment for shockable cardiac arrest rhythms (VF/pulseless VT).

\section*{How This Test or Procedure Is Performed}
Defibrillator pads are placed on the patient's chest (anterior-lateral or anterior-posterior). The machine is charged, a 'clear' command is given, and an electrical shock is delivered.

\section*{What Preparation Is Needed}
Apply conductive gel or adhesive pads to the chest. Ensure the patient is not in standing water and no one is touching the patient.

\section*{Contraindications and Precautions}
Asystole and Pulseless Electrical Activity (PEA) are non-shockable rhythms; defibrillation is contraindicated and CPR should continue.

\section*{Risks}
Skin burns at the pad sites, skeletal muscle injury, or accidental shock to healthcare providers.

\section*{What Happens After the Test or Procedure}
Immediately resume chest compressions for 2 minutes before re-analyzing the heart rhythm.

\section*{Understanding Your Results}
The rhythm is analyzed after 2 minutes of CPR to check if the shockable rhythm was terminated and replaced by an organized rhythm.

\section*{Normal Results}
Termination of VF/VT; restoration of a perfusing sinus rhythm.

\section*{Follow-up Tests and Procedures}
Antiarrhythmic medications (amiodarone), post-arrest care, or implantable cardioverter-defibrillator (ICD) placement.

\section*{References}
\begin{enumerate}
  \item MedlinePlus. Defibrillation. U.S. National Library of Medicine; 2025.
  \item Current Medical Diagnosis and Treatment. McGraw-Hill; 2025.
\end{enumerate}

\clearpage
